%------------------------------------------------------------------------------%
\begin{question}
  Determine a posição relativa entre as retas
  \[
    r\colon
    \begin{cases}
      x = 2+t   \\
      y = -1+2t \\
      z = 1-t
    \end{cases}
  \qquad\qquad
    s\colon
    \begin{cases}
      x = 1+2\tau \\
      y = 1+4\tau \\
      z = 3-2\tau
    \end{cases}
  \]
\end{question}

%------------------------------------------------------------------------------%
\begin{resolution}{Solução}
  Os vetores diretores são
  \[
    \vec{v} = (1, 2, -1)
  \]
  \[
    \vec{w} = (2, 4, -2)
  \]
  \[
    \vec{w} = 2\vec{v}
  \]
  as retas são paralelas ou coincidentes
\end{resolution}

%------------------------------------------------------------------------------%
\begin{resolution}{Solução}
  Escolhemos
  \[
  P_0=(2,-1,1)
  \]
  e
  \[
  Q_0=(1,1,3)
  \]
  Então,
  \[
  Q_0-P_0=(-1,2,2)
  \]
  Se as retas fossem coincidentes, esse vetor deveria ser paralelo a
  \[
  \vec v=(1,2,-1)
  \]
  Isso não ocorre.

  De fato, comparando as componentes:
  \[
  \frac{-1}{1}\neq\frac{2}{2}
  \]
  Logo, as retas são \textbf{paralelas distintas}.
\end{resolution}

%------------------------------------------------------------------------------%
\endinput